%! Author = sriadityavedantam
%! Date = 3/30/26

\section{Chapter 6}
\begin{problem}{
    Let $\mathbb{F}$ be a field and $f(x), g(x), p(x)\in\mathbb{F}[x]$ with $p(x)$ nonzero.
    Then $f(x)$ is congruent to $g(x)$ modulo $p(x)$, written
    \[
        f(x)\equiv g(x)\mod p(x),
    \]
    provided that $p(x)$ divides $f(x)-g(x)$.
    Prove or disprove:
    If $p(x)$ is relatively prime to $k(x)$ and
    \[
        f(x)k(x)\equiv g(x)k(x)\mod p(x),
    \]
    then
    \[
        f(x)\equiv g(x)\mod p(x).
    \]
}
    This statement is true.
    Suppose
    \[
        \gcd(p(x),k(x))=1_{\mathbb{F}}
    \]
    and
    \[
        f(x)k(x)\equiv g(x)k(x)\mod p(x).
    \]
    By the definition of polynomial congruence modulo $p(x)$, this means
    \[
        p(x)\mid f(x)k(x)-g(x)k(x).
    \]
    Thus
    \begin{align*}
        p(x) &\mid k(x)(f(x)-g(x)).
    \end{align*}
    Since $p(x)$ is relatively prime to $k(x)$, Euclid's Lemma for polynomials implies that
    \[
        p(x)\mid f(x)-g(x).
    \]
    Therefore
    \[
        f(x)\equiv g(x)\mod p(x).
    \]
    So the statement is proved.
\end{problem}
\begin{problem}{
    Let $\mathbb{F}$ be a field and $f(x), g(x), p(x)\in\mathbb{F}[x]$ with $p(x)$ nonzero.
    Then $f(x)$ is congruent to $g(x)$ modulo $p(x)$, written
    \[
        f(x)\equiv g(x)\mod p(x),
    \]
    provided that $p(x)$ divides $f(x)-g(x)$.
    Prove or disprove:
    If $p(x)$ is irreducible in $\mathbb{F}[x]$ and
    \[
        f(x)g(x)\equiv 0\mod p(x),
    \]
    then
    \[
        f(x)\equiv 0\mod p(x)
        \quad\text{or}\quad
        g(x)\equiv 0\mod p(x).
    \]
}
    This statement is true.
    Suppose $p(x)$ is irreducible in $\mathbb{F}[x]$ and
    \[
        f(x)g(x)\equiv 0\mod p(x).
    \]
    By the definition of polynomial congruence modulo $p(x)$, this means
    \[
        p(x)\mid f(x)g(x).
    \]
    Since $p(x)$ is irreducible, Euclid's Lemma for polynomials implies that
    \[
        p(x)\mid f(x)
        \quad\text{or}\quad
        p(x)\mid g(x).
    \]
    This is equivalent to
    \[
        p(x)\mid f(x)-0
        \quad\text{or}\quad
        p(x)\mid g(x)-0.
    \]
    Hence
    \[
        f(x)\equiv 0\mod p(x)
        \quad\text{or}\quad
        g(x)\equiv 0\mod p(x).
    \]
\end{problem}
\begin{problem}{
    Write out the addition and multiplication tables for the congruence class ring
    \[
        \mathbb{Z}_3[x]/(x^2+1).
    \]
    Is $\mathbb{Z}_3[x]/(x^2+1)$ a field?
}
    \[
        x^2\equiv -1\equiv 2 \mod (x^2+1),
    \]
    so every congruence class has a representative of the form
    \[
        a+bx
    \]
    with $a,b\in\mathbb{Z}_3$.
    Thus the elements are
    \[
        0,\ 1,\ 2,\ x,\ x+1,\ x+2,\ 2x,\ 2x+1,\ 2x+2.
    \]
    \begin{center}
        \begin{tabular}{|c|c|c|c|c|c|c|c|c|c|}
            \hline
            $+$ & $0$ & $1$ & $2$ & $x$ & $x+1$ & $x+2$ & $2x$ & $2x+1$ & $2x+2$ \\ \hline
            $0$ & $0$ & $1$ & $2$ & $x$ & $x+1$ & $x+2$ & $2x$ & $2x+1$ & $2x+2$ \\ \hline
            $1$ & $1$ & $2$ & $0$ & $x+1$ & $x+2$ & $x$ & $2x+1$ & $2x+2$ & $2x$ \\ \hline
            $2$ & $2$ & $0$ & $1$ & $x+2$ & $x$ & $x+1$ & $2x+2$ & $2x$ & $2x+1$ \\ \hline
            $x$ & $x$ & $x+1$ & $x+2$ & $2x$ & $2x+1$ & $2x+2$ & $0$ & $1$ & $2$ \\ \hline
            $x+1$ & $x+1$ & $x+2$ & $x$ & $2x+1$ & $2x+2$ & $2x$ & $1$ & $2$ & $0$ \\ \hline
            $x+2$ & $x+2$ & $x$ & $x+1$ & $2x+2$ & $2x$ & $2x+1$ & $2$ & $0$ & $1$ \\ \hline
            $2x$ & $2x$ & $2x+1$ & $2x+2$ & $0$ & $1$ & $2$ & $x$ & $x+1$ & $x+2$ \\ \hline
            $2x+1$ & $2x+1$ & $2x+2$ & $2x$ & $1$ & $2$ & $0$ & $x+1$ & $x+2$ & $x$ \\ \hline
            $2x+2$ & $2x+2$ & $2x$ & $2x+1$ & $2$ & $0$ & $1$ & $x+2$ & $x$ & $x+1$ \\ \hline
        \end{tabular}
    \end{center}
    \begin{center}
        \begin{tabular}{|c|c|c|c|c|c|c|c|c|c|}
            \hline
            $\cdot$ & $0$ & $1$ & $2$ & $x$ & $x+1$ & $x+2$ & $2x$ & $2x+1$ & $2x+2$ \\ \hline
            $0$ & $0$ & $0$ & $0$ & $0$ & $0$ & $0$ & $0$ & $0$ & $0$ \\ \hline
            $1$ & $0$ & $1$ & $2$ & $x$ & $x+1$ & $x+2$ & $2x$ & $2x+1$ & $2x+2$ \\ \hline
            $2$ & $0$ & $2$ & $1$ & $2x$ & $2x+2$ & $2x+1$ & $x$ & $x+2$ & $x+1$ \\ \hline
            $x$ & $0$ & $x$ & $2x$ & $2$ & $x+2$ & $2x+2$ & $1$ & $x+1$ & $2x+1$ \\ \hline
            $x+1$ & $0$ & $x+1$ & $2x+2$ & $x+2$ & $2x$ & $1$ & $2x+1$ & $2$ & $x$ \\ \hline
            $x+2$ & $0$ & $x+2$ & $2x+1$ & $2x+2$ & $1$ & $x$ & $x+1$ & $2x$ & $2$ \\ \hline
            $2x$ & $0$ & $2x$ & $x$ & $1$ & $2x+1$ & $x+1$ & $2$ & $2x+2$ & $x+2$ \\ \hline
            $2x+1$ & $0$ & $2x+1$ & $x+2$ & $x+1$ & $2$ & $2x$ & $2x+2$ & $x$ & $1$ \\ \hline
            $2x+2$ & $0$ & $2x+2$ & $x+1$ & $2x+1$ & $x$ & $2$ & $x+2$ & $1$ & $2x$ \\ \hline
        \end{tabular}
    \end{center}
    Yes, $\mathbb{Z}_3[x]/(x^2+1)$ is a field.
    The reason is that $x^2+1$ is irreducible in $\mathbb{Z}_3[x]$.
    Indeed, the only elements of $\mathbb{Z}_3$ are $0,1,2$, and
    \[
        0^2+1=1,\qquad 1^2+1=2,\qquad 2^2+1=4+1\equiv 2 \mod 3,
    \]
    so $x^2+1$ has no root in $\mathbb{Z}_3$.
    Hence it is irreducible, and the quotient is a field.
\end{problem}
\begin{problem}{
    Let $F$ be a field.
    A nonconstant polynomial $p(x)\in F[x]$ is said to be irreducible if its only divisors are its associates and the nonzero constant polynomials.
    In each part, explain why $[f(x)]$ is a unit in $\mathbb{F}[x]/(p(x))$ and find its inverse.

    (a) $[f(x)] = [2x-3]\in\mathbb{Q}[x]/(x^2-2)$.

    (b) $[f(x)] = [x^2+x+1]\in\mathbb{Z}_3[x]/(x^2+1)$.
}
    \textbf{(a).}
    Since $x^2-2$ is irreducible in $\mathbb{Q}[x]$, the quotient
    \[
        \mathbb{Q}[x]/(x^2-2)
    \]
    is a field.
    Thus every nonzero class is a unit.
    Now $[2x-3]\neq [0]$ in this quotient, since otherwise $x^2-2$ would divide the linear polynomial $2x-3$, which is impossible.
    So $[2x-3]$ is a unit.
    To find its inverse, we solve
    \[
        (2x-3)(ax+b)\equiv 1\mod (x^2-2).
    \]
    Using $x^2\equiv 2$, we compute
    \begin{align*}
        (2x-3)(ax+b)
&=
2ax^2+(2b-3a)x-3b\\
&\equiv
4a+(2b-3a)x-3b.
    \end{align*}
    So we want
    \[
        (2b-3a)x+(4a-3b)=1.
    \]
    Thus
    \begin{align*}
        2b-3a &= 0,\\
        4a-3b &= 1.
    \end{align*}
    From the first equation, $2b=3a$.
    Over $\mathbb{Q}$ this gives
    \[
        b=\frac{3a}{2}.
    \]
    Substitute into the second:
    \[
        4a-3\left(\frac{3a}{2}\right)=1,
    \]
    so
    \[
        \frac{8a-9a}{2}=1,
    \]
    hence
    \[
        -\frac{a}{2}=1,
        \qquad
        a=-2,
        \qquad
        b=-3.
    \]
    Therefore
    \[
        [2x-3]^{-1}=[-2x-3].
    \]

    \textbf{(b).}
    In $\mathbb{Z}_3[x]/(x^2+1)$, we first reduce
    \[
        x^2+x+1 \equiv 2+x+1 = x
    \]
    since $x^2\equiv -1\equiv 2$ modulo $x^2+1$.
    So
    \[
        [x^2+x+1]=[x].
    \]
    Since the quotient is a field, any nonzero element is a unit.
    Also $[x]\neq [0]$, so $[x^2+x+1]$ is a unit.
    To find its inverse, note that
    \[
        x^2\equiv 2 \mod (x^2+1),
    \]
    so
    \[
        x\cdot 2x = 2x^2 \equiv 2\cdot 2 = 4 \equiv 1 \mod 3.
    \]
    Therefore
    \[
        [x]^{-1}=[2x].
    \]
    Hence
    \[
        [x^2+x+1]^{-1}=[2x].
    \]
\end{problem}

\begin{problem}{
    Find a fourth-degree polynomial in $\mathbb{Z}_2[x]$ whose roots are the four elements of the field
    \[
        \mathbb{Z}_2[x]/(x^2+x+1).
    \]
}
    The field
    \[
        \mathbb{Z}_2[x]/(x^2+x+1)
    \]
    has four elements:
    \[
        [0],\ [1],\ [x],\ [x+1].
    \]
    A polynomial whose roots are exactly these four elements is
    \[
        f(t)=t(t+1)(t+x)(t+x+1).
    \]
    Over a field with $4$ elements, every element $a$ satisfies
    \[
        a^4=a,
    \]
    so the polynomial
    \[
        t^4-t
    \]
    has all four field elements as roots.
    Since we are in characteristic $2$, this is the same as
    \[
        t^4+t.
    \]
    Thus one such fourth-degree polynomial is
    \[
        f(t)=t^4+t.
    \]
    Indeed,
    \[
        f(0)=0,
        \qquad
        f(1)=1+1=0,
        \qquad
        f(x)=x^4+x=0,
        \qquad
        f(x+1)=(x+1)^4+(x+1)=0.
    \]
    So the four elements of the field are precisely the roots.
\end{problem}

\begin{problem}{
    (a) Show that $\mathbb{Z}_2[x]/(x^3+x+1)$ is a field.

    (b) Show that the field $\mathbb{Z}_2[x]/(x^3+x+1)$ contains all three roots of $x^3+x+1$.
}
    \textbf{(a).}
    We show that
    \[
        x^3+x+1
    \]
    is irreducible in $\mathbb{Z}_2[x]$.
    A cubic over a field is reducible if and only if it has a root.
    So check the two elements of $\mathbb{Z}_2$:
    \begin{align*}
        f(0) &= 0^3+0+1 = 1 \neq 0,\\
        f(1) &= 1^3+1+1 = 1 \neq 0
    \end{align*}
    in $\mathbb{Z}_2$.
    Thus $x^3+x+1$ has no root in $\mathbb{Z}_2$, so it is irreducible.
    Therefore
    \[
        \mathbb{Z}_2[x]/(x^3+x+1)
    \]
    is a field.

    \textbf{(b).}
    Let
    \[
        \alpha=[x]\in \mathbb{Z}_2[x]/(x^3+x+1).
    \]
    Then
    \[
        \alpha^3+\alpha+1=0
    \]
    in the quotient, since
    \[
        x^3+x+1\equiv 0 \mod (x^3+x+1).
    \]
    So $\alpha$ is one root.
    Since we are in characteristic $2$, the nonzero elements of this field form a multiplicative group of order $7$.
    Thus every nonzero element satisfies $u^7=1$.
    Now from
    \[
        \alpha^3+\alpha+1=0
    \]
    we get
    \[
        \alpha^3=\alpha+1.
    \]
    Then
    \begin{align*}
        (\alpha^2)^3+\alpha^2+1
    &=
        \alpha^6+\alpha^2+1.
    \end{align*}
    Because $\alpha^7=1$, we have
    \[
        \alpha^6=\alpha^{-1}.
    \]
    Also, from $\alpha^3=\alpha+1$, multiplying by $\alpha^{-1}$ gives
    \[
        \alpha^2=1+\alpha^{-1},
    \]
    so
    \[
        \alpha^{-1}=\alpha^2+1.
    \]
    Hence
    \[
        \alpha^6+\alpha^2+1 = (\alpha^2+1)+\alpha^2+1=0.
    \]
    So $\alpha^2$ is also a root.
    Similarly,
    \begin{align*}
        (\alpha^4)^3+\alpha^4+1
    &=
        \alpha^{12}+\alpha^4+1\\
        &=
        \alpha^5+\alpha^4+1
    \end{align*}
    since $\alpha^7=1$.
    Now
    \[
        \alpha^4=\alpha(\alpha^3)=\alpha(\alpha+1)=\alpha^2+\alpha,
    \]
    and
    \[
        \alpha^5=\alpha(\alpha^4)=\alpha(\alpha^2+\alpha)=\alpha^3+\alpha^2=(\alpha+1)+\alpha^2.
    \]
    Thus
    \[
        \alpha^5+\alpha^4+1=(\alpha^2+\alpha+1)+(\alpha^2+\alpha)+1=0.
    \]
    So $\alpha^4$ is also a root.
    Therefore the three roots of $x^3+x+1$ in this field are
    \[
        [x],\ [x^2],\ [x^4].
    \]
    Hence the field contains all three roots.
\end{problem}
\begin{problem}{
    Show that every polynomial of degree $1$, $2$, or $4$ in $\mathbb{Z}_2[x]$ has a root in
    \[
        \mathbb{Z}_2[x]/(x^4+x+1).
    \]
}
    Let
    \[
        E\coloneqq \mathbb{Z}_2[x]/(x^4+x+1).
    \]
    Since $x^4+x+1$ is irreducible over $\mathbb{Z}_2$, $E$ is a field with
    \[
        2^4=16
    \]
    elements.
    A finite field of order $16$ has multiplicative group of order $15$, so every nonzero element $a\in E$ satisfies
    \[
        a^{15}=1.
    \]
    Thus every nonzero element is a root of
    \[
        t^{15}-1.
    \]
    Since we are in characteristic $2$, this is
    \[
        t^{15}+1.
    \]
    So all $16$ elements of $E$ are roots of
    \[
        t^{16}-t.
    \]
    Now let $f(t)\in\mathbb{Z}_2[t]$ be irreducible of degree $1$, $2$, or $4$.
    Then the splitting field of $f$ over $\mathbb{Z}_2$ has size $2^d$, where $d=\deg f$.
    Since $d$ is one of $1,2,4$, we have
    \[
        d\mid 4.
    \]
    Hence the field with $2^d$ elements embeds into the field with $2^4=16$ elements.
    Therefore every irreducible polynomial of degree $1$, $2$, or $4$ over $\mathbb{Z}_2$ splits in $E$, and in particular has a root in $E$.
    Since every polynomial of degree $1$, $2$, or $4$ factors into irreducibles whose degrees are among $1$, $2$, and $4$, it follows that every such polynomial has a root in $E$.
    Thus every polynomial of degree $1$, $2$, or $4$ in $\mathbb{Z}_2[x]$ has a root in
    \[
        \mathbb{Z}_2[x]/(x^4+x+1).
    \]
\end{problem}